% !TEX root = main.tex
\section{Experimental Setup}
\label{sec:eval:setup}

\myparagraph{Models and tasks}
Our primary workload is AG News four-class topic classification,
partitioned non-IID across clients by a Dirichlet distribution
($\alpha = 1$). \tbd{Extend to two further datasets --- Yahoo! Answers
topic classification and D3 --- and two further models beyond
DistilBERT (M2, M3), to show the gain is not an artifact of one
model/task pairing; not yet run.} We fine-tune DistilBERT-base-uncased
(66.4\,M parameters, backbone frozen) through AdapterHub bottleneck
adapters~\cite{houlsby2019parameter,pfeiffer2020adapterhub}
(\tbd{$\sim$1.5\%} of parameters trainable). This is a
parameter-efficient configuration whose peak training memory is
bounded by forward inference rather than backpropagation, so it maps
onto inference-only mobile accelerators by construction --- the
mobile-feasibility argument is structural, not measured on-device (see
\S\ref{sec:bg}). Generalization to larger model families (LLaMA2-7B,
Mistral-7B) and to more severe non-IID is deferred to the ablations
(\S\ref{sec:ablation}). \tbd{Bridge to the memory-wall motivation: the
argument that forces forward gradients (\S\ref{sec:bg}) is a 7B-scale
claim, whereas this testbed is DistilBERT; state why the smaller,
tractable testbed does not undercut the 7B motivation before
submission.}

\myparagraph{Testbed and modeled clock}
Each trainer and the aggregator runs as a separate process on an
8$\times$NVIDIA~A40 server, exchanging real messages over MQTT;
on-device mobile latency is \emph{modeled}, not measured. Each
forward-only round trip is drawn from the FedBuff/Papaya mobile-device
runtime distribution~\cite{nguyen2022fedbuff}\tbd{add Papaya cite} and
scaled down by \texttt{delay\_factor} (\tbd{=2}) to reflect that
forward-only inference on an NPU costs a fraction of a backward pass.
\dhruv{But this was based on runtimes and forward pass estimates per
baseline and extrapolated to the Papaya distribution.} \dhruv{verify
from code what factor we use finally.} We evaluate two availability
conditions: 100\% (algorithmic isolation; no dropout) and a
\tbd{mobiperf} device-availability trace as the real-world deployment
condition. We deliberately do \emph{not} re-run a synthetic
availability sweep here: the motivation section
(\S\ref{sec:motiv:control}) runs that sweep pedagogically. Results are
reported for a single seed \tbd{(median over $\ge$3 seeds for the
headline time-to-accuracy is a TODO before submission)}. \dhruv{Should
list out what all was seeded at the aggregator and trainer so as to
help us make the case that the results are reproducible.}

\myparagraph{Baselines}
\tabref{tab:eval:baselines} lists the nine configurations compared
against \sys{} and defines their naming grammar in the caption.
\textbf{+O} is inert under 100\% availability, since there is nothing
to track --- it diverges from \textbf{+IT} only under the real-world
trace. Four rows --- bold, above the rule in
\tabref{tab:eval:baselines} --- are published anchors compared against
\sys{} directly: \fwdllm{}, the only prior forward-gradient FL system;
\fedbuffp{}~\cite{nguyen2022fedbuff} and
\felixp{}~\cite{garg2026felix}, whose selection and aggregation
strategies are ported from published async-FL systems onto our
substrate (exact knobs in the table); and \sys{} itself. The five rows
below the rule are our own intermediate configurations that isolate
one axis of change at a time and are never presented as a headline
comparison. Prior work documents the same availability-unawareness gap
this baseline set is designed to probe~\cite{garg2025avail}.

To keep the nine rows easy to place, each varies exactly one of three
things relative to \fwdllm{}: \emph{how} clients synchronize (sync
vs.\ async), \emph{when} control decisions are made (round vs.\
iteration), and \emph{which} clients are selected (uniform vs.\
availability-aware substrate). The three axes below walk through these
in turn. (1)~\emph{Sync $\to$ async}: \fwdllm{} $\to$ \fedbuffp{}
holds selection and control granularity fixed and removes only the
sync barrier, isolating whether asynchronous execution alone helps
before any contribution of ours. (2)~\emph{Round $\to$ iteration},
checked independently on three selection substrates --- \fwdllm{}
$\to$ \fwdllmit{}, \fedbuffp{} $\to$ \fedbuffpit{}, and \felixp{}
$\to$ \felixpit{} --- isolates the iteration-level control reframe and
confirms the gain is not an artifact of one particular selector.
\fedbuffpit{} is the direct attribution floor for \sys{}'s
contributions: it is \sys{} with C1, C2, and C3 removed and everything
else held fixed, so the residual gap to \sys{} is exactly what the
three contributions buy. (3)~\emph{Selection substrate}: \felixp{}
asks whether availability-aware smart selection alone, without any of
\sys{}'s own contributions, already closes most of the gap;
\felixpit{} transplants \sys{}'s iteration-level cadence onto FeLiX's
selector to isolate that cadence from C1/C3 specifically, and is never
presented as a headline comparison, since crediting a published
baseline with our own mechanism would be an unearned advantage to that
baseline. \felixp{} and \felixpit{} are \emph{our} re-run of FeLiX's
ported strategy on this paper's task, never FeLiX's own reported
numbers --- FeLiX's evaluation is image/speech classification on a
disjoint substrate~\cite{garg2026felix}, and results never cross
between the two papers in either direction.

\input{tables/eval/baselines}

All baselines run at their best-tuned hyperparameters; the tuning
procedure is in the table \tbd{add tuning-procedure table}. \fwdllm{}
runs in its native round-level configuration, not retrofitted with
iteration-level structure. SPRY~\cite{panchal2024spry} is discussed as
a gradient-quality upper bound but is not a run-for-run baseline. Like
other methods that \emph{split or personalize the model per device}
\tbd{(cite split-learning e.g.\ Vepakomma'18; model-heterogeneous FL
e.g.\ HeteroFL/Diao'21; personalized-FL survey e.g.\ Tan'22)}, its
per-client weight splitting is incompatible with the homogeneous
cross-device deployment we target --- a fleet of interchangeable,
intermittently-available clients that each run the \emph{same} model
and can drop in and out, where a device-specific model shard cannot
survive churn. We justify this exclusion where these methods are
introduced (\S\tbd{background/related work}); here it only scopes the
baseline set.

\myparagraph{Excluded selection/staleness baselines}
Three further classical FL baselines are cited as ancestors rather
than run for run, and none change the scope of
\tabref{tab:eval:baselines}. Oort's~\cite{oort} multi-armed-bandit
utility-guided selection is already the idea inside \felixp{}'s (and
\sys{}'s own) selector, so re-running it separately would be
redundant. REFL's~\cite{abdelmoniem2023refl} deadline/overcommit
staleness model does not map onto a variance-gated forward-gradient
commit, since our clients take no local update steps for a deadline to
bound. FedDance's availability-forecasting selection targets a
trace-driven evaluation we defer to future work.

\tbd{\fedbuffp{}, \fedbuffpit{}, \fedbuffpito{}, \felixp{}, and
\felixpit{} are specified above but not yet wired into the run
pipeline; \S\ref{sec:eval:sota} and \S\ref{sec:eval:attribution}
currently report only the \fwdllm{}/\fwdllmito{}/\sys{} rung that has
completed runs, and will add the remaining anchors' numbers before
submission.}

\myparagraph{Baseline fidelity and the simulation clock}
We port \fwdllm{} onto Flame~\cite{daga2023flame} so that it and
\sys{} run on the same faithful platform, and we confirm the port is
fidelity-preserving by \emph{accuracy parity}: our port reaches the
test accuracy the original paper reports~\cite{xu2024fwdllm}
\tbd{state the matched value, e.g.\ ours [X]\% vs.\ their [Y]\%}. We
compare accuracy, not wall-clock, and we say so deliberately: a
wall-clock comparison against the original is not meaningful, because
that system lacks the real distributed implementation ours reproduces
(separate trainer processes, real network messaging, per-trainer
state). For the same reason, the two properties below make our
platform harder on the baseline, not easier. First, Flame runs each
trainer and the aggregator as separate processes exchanging real
messages over MQTT, whereas the original ran a single process with one
aggregator and one client emulating all trainers in a closed loop.
Faithful distribution forced us to maintain per-trainer state that the
single-process baseline read for free from its parent, but that a real
deployment would not have. Second, per-client round-trip times are
drawn from the modeled mobile-delay distribution described above,
scaled by \texttt{delay\_factor} --- accelerated inference on an NPU
costs a fraction of the backward pass those traces measured. This
modeled clock, not a validated virtual clock, is what our
time-to-accuracy uses; a faithful real$\leftrightarrow$sim virtual
clock is future work. Every reported speedup is therefore against a
correctly distributed, correctly aggregated baseline.

% =======================================================================
\section{Evaluation}
\label{sec:eval}

Time-to-accuracy is the metric \sys{} optimizes; resource utilization,
compute effectiveness, communication, and session length are not
separate experiments but four further readings off the \emph{same} set
of executions. We run training once per baseline per availability
condition, and every number below is a query over those logs.
\S\ref{sec:eval:sota} answers two questions --- \emph{how much faster}
does \sys{} reach target accuracy than each published anchor
(\fwdllm{}, \fedbuffp{}, \felixp{}), and \emph{how cheap} are those
wins for the device, in trainer idling, data moved, and session length
--- leading with \fwdllm{} since it is the only prior
\emph{forward-gradient} FL system, so the size of that gap is the
reader's first impression, and situating \fedbuffp{} and \felixp{}
alongside it as the async-execution and smart-selection anchors.
\S\ref{sec:eval:attribution} then asks where the advantage comes from:
against baselines we \emph{strengthened} with the obvious systems
fixes (the +IT/+O rungs in \tabref{tab:eval:baselines}) and the 100\%
oracle, do \sys{}'s \emph{conceptual} contributions still pay off, and
is the residual gap explained by the two mechanistic levers --- higher
utilization (less trainer/aggregator wait) and higher compute
effectiveness (more learning per unit compute)?

% =======================================================================
\subsection{How much does \sys improve end-to-end LLM fine-tuning in FL?}
\label{sec:eval:sota}

\fwdllm{} is the sole prior forward-gradient FL system, so we report
the absolute gap to it first, across all five metrics on one run-set;
\fedbuffp{} and \felixp{} extend the same comparison along the
async-execution and smart-selection axes and will be added once their
runs complete (\S\ref{sec:eval:setup}). Under the real-world
availability trace, \fwdllm{}'s round commitment collides with short
availability windows and it barely learns; the gaps below are
correspondingly large.

\input{figs/code/eval/e2e/e1_tta}

\myparagraph{Time to accuracy}
\textbf{\sys{} reaches the accuracy target while \fwdllm{} never
converges.} On AG News ($N{=}100$, non-IID $\alpha{=}1$), \sys{}
crosses the $84\%$ target in $3.9$\,h, whereas \fwdllm{} plateaus at
$30.3\%$ --- its whole-round commitment discards the in-flight
databins of every client that drops inside a round, so it never
accumulates enough usable updates to learn (\figref{fig:eval:tta}).
\label{sec:eval:tta} This is the outcome that matters: round-level
control turns the real-world availability scenario into a
near-failure, and iteration-level control turns it back into a
tractable run. We report the convergence event in four units ---
wall-clock, rounds, data bins, and iterations --- plus the maximum
accuracy attained; virtual-clock time is deferred (real-mode runs emit
no virtual clock and the sim clock is not yet validated).

\myparagraph{Resource utilization}
\textbf{\sys{} does not ask clients for more compute --- it wastes
less of the compute each client already offers.} Every client's budget
is set by the same thing regardless of scheduler: the wall-clock
window it happens to be reachable for. A round-barrier scheduler
cannot use that window until the round it belongs to closes, so the
client sits idle mid-window waiting on its slower peers; \sys{}'s
iteration-level reselection re-dispatches a client the moment it
returns a scalar, so the same window is filled with work instead of
wait. Decomposing each run's timeline into busy and idle fractions
makes this concrete: \sys{}'s per-trainer busy fraction is higher at
every percentile of the fleet (median ${\sim}7$--$8\%$), not because
\sys{} demands more from clients, but because the baseline leaves that
fraction of each client's own availability window unused
(\figref{fig:eval:util}). \label{sec:eval:util} The aggregator stays
compute-bound rather than wait-bound, confirming the coordinator
itself is not what is idling. Higher utilization is therefore a
genuine systems win, not a hidden compute subsidy: it is compute the
client's own availability window already afforded, that only \sys{}'s
control structure knows how to claim.

\input{figs/code/eval/e2e/e2_util}

\myparagraph{Compute effectiveness}
\textbf{\sys{} converts client compute into learning where \fwdllm{}
regresses.} Against a hardware-independent forward-pass denominator,
\sys{} yields $+1.68$ $\Delta\text{loss}$ per million forward passes
while \fwdllm{} \emph{loses} ground ($-1.26$): \fwdllm{}'s random
directions overwhelmingly return near-zero scalars, so the compute its
clients spend buys essentially no learning
(\figref{fig:eval:compute}). \label{sec:eval:compute} We lead with the
forward-pass denominator because \texttt{gpu\_compute\_s} is wall-time
under 100-trainer / 8-GPU contention, which inflates it
(${\sim}20\times$) and differently across baselines, so it measures
scheduling contention rather than algorithmic compute; we report the
GPU-hour denominator as a confounded secondary. Guided perturbations
are the effectiveness lever behind this gap: a client that evaluates
an informative direction reduces loss more than one that evaluates a
random direction at the same cost.

\input{figs/code/eval/e2e/e3_compute}

\myparagraph{Communication}
\textbf{Per-message uplink is a scalar, so total bytes stay modest,
but \sys{}'s continuous re-pulls raise its volume in this C1-only
configuration.} We measure the data actually moved in both directions
(\figref{fig:eval:comm}). \label{sec:eval:comm} \sys{} sends more,
smaller messages per round; because the uplink is scalar-dominated,
the message-count increase does not by itself inflate bytes. In the
current C1-only run-set, however, \sys{}'s async clients continuously
re-pull the global model, so its total uplink volume sits \emph{above}
the synchronous baselines rather than below --- the honest current
state (\S\ref{sec:eval:attribution} treats this directly). \weak{The
intended net reduction depends on C2 (dynamic $K$/$C$, fewer
concurrent re-pulls) and delta-model distribution, not yet enabled;
until then, report the measured bytes honestly.}

\input{figs/code/eval/e2e/e4_comm}

\myparagraph{Client session length}
\dhruv{it seems that the telemetry for this is not being captured correctly. the session lengths for baselines are the continuous wall clock time for which a trainer remains selected in the cohort. It is not that intermittent training time from the trainer's perspective.}
\textbf{\sys{}'s client sessions are short, so intermittent devices
can participate.} A round-based system needs a client reachable for a
whole round; \sys{} needs it for one iteration. Measured
active-session durations are short (seconds-scale) across all three
systems, but \sys{}'s concentrate in a tight band that matches its
per-iteration engagement model (\figref{fig:eval:sessions}).
\label{sec:eval:sessions} This is the mechanism behind \sys{}'s
robustness to the short availability windows documented in
\S\ref{sec:motiv:control}: a client need only stay reachable for
seconds-to-minutes, not for a full round.

\input{figs/code/eval/e2e/e5_sessions}

% =======================================================================
\subsection{Where does the advantage come from?}
\label{sec:eval:attribution}

The gap to \fwdllm{} is large but conflates two things: the systems
fixes any competent engineer would add (per-iteration reselection,
oracular availability) and \sys{}'s \emph{novel} mechanisms (guided
perturbations, dynamic $K$/$C$, gradient-aware aggregation). To
isolate the second, we compare against \fwdllmito{} --- \fwdllm{}
\emph{with} those systems fixes --- and the 100\% oracle.
\tbd{\fedbuffpito{} and \felixpit{} extend the same attribution once
their runs complete, holding the strengthened-baseline standard to the
async-execution and smart-selection anchors too.} \fwdllmito{} is the
iteration-level reframe \emph{without} our contributions; that it
still trails \sys{} sharply is the evidence that the reframe is an
enabler, not the gain --- moving control to the iteration opens the
decision space, but the contributions are what exploit it.
\fwdllmito{} is expected to beat \fwdllm{} by construction, so the
remaining gap to \sys{} is smaller; the point of this subsection is
not the size of that gap but its \emph{source}. We lead with
time-to-accuracy and the two levers that explain it, and we state
honestly where the current C1-only run-set does not yet realize a
claimed win.

\myparagraph{Time to accuracy vs.\ the strengthened baseline}
\textbf{Even after the systems fixes, \sys{} is ${\sim}2\times$
faster.} \fwdllmito{} peaks at $80.9\%$ and takes $7.4$\,h --- about
$2\times$ \sys{}'s $3.9$\,h --- and does not reach the $84\%$ target
(\figref{fig:eval:tta}). Per-iteration reselection and oracular
tracking recover most of \fwdllm{}'s collapse, but the residual gap is
what the novel mechanisms buy: the strengthened baseline still commits
to whole-databin work and static $K$/$C$, so it converges more slowly
and to a lower peak.

\myparagraph{Utilization vs.\ the oracular scheduler}
\textbf{Oracular availability does not close the utilization gap.}
\fwdllmito{}'s oracular sync selection still leaves most clients
near-idle (${\sim}1\%$ busy) because its synchronization barrier
stalls fast clients waiting on slow ones, whereas \sys{}'s
barrier-free reselection keeps them busy (\figref{fig:eval:util}).
Knowing device state perfectly does not help if the control structure
cannot act on it per iteration --- attribution the oracle isolates
cleanly, since it removes availability uncertainty as a confound.

\myparagraph{Compute effectiveness: an honest caveat}
\textbf{Per unit compute, the strengthened baseline is currently more
efficient than \sys{}.} On the forward-pass denominator, \fwdllmito{}
reaches $4.58$ $\Delta\text{loss}$/M-fwd against \sys{}'s $1.68$
(\figref{fig:eval:compute}). This is expected and disclosed: guidance
pays a $2P$ per-iteration surcharge (\S\ref{sec:arch:c1}), and \sys{}
trades per-compute efficiency for wall-clock speed with its C2/C3
efficiency levers \emph{not yet enabled} in this run-set. \sys{} still
wins wall-clock time-to-accuracy because it needs far fewer
iterations; the per-unit efficiency win is the hypothesis that the
efficiency levers close. \weak{The surcharge is repaid per unit
compute only if C2 (dynamic $K$/$C$) and C3 (gradient-aware
aggregation) cut iterations by more than the ${\sim}10\times$
per-iteration factor at $P{=}10$. On the current C1-only run-set they
are not yet enabled, so the oracular sync baseline is more per-compute
efficient. Report the forward-pass denominator honestly and scope the
compute claim to the constrained-client regime until C2/C3 land.}

\myparagraph{Communication: an honest caveat}
\textbf{\sys{}'s total bytes are currently higher than the oracular
baseline, not lower.} As above, \sys{}'s async clients re-pull the
global model continuously, so on the C1-only run-set total uplink is
higher than \fwdllmito{} (\figref{fig:eval:comm}). The reduction is
the hypothesis that C2 (fewer concurrent re-pulls) and fewer
iterations-to-target deliver; on the full-system run we expect the net
effect to invert. Until then we report the measured bytes as they are.

\myparagraph{Accounting for the guidance surcharge}
Guided selection raises \sys{}'s per-iteration compute
(\S\ref{sec:arch:c1}), so we report the surcharge openly. At our
default $P=10$ candidate perturbations, \sys{} evaluates $2P=20$
forward passes per iteration against the sync baseline's $2$ --- about
$10\times$ the per-iteration compute. The fidelity-preserving
optimizations of \S\ref{sec:arch:c1} are bit-identical and cut
\sys{}'s per-iteration GPU time by ${\sim}37\%$ (removing 5 redundant
passes, $25\to20$); they lower the absolute cost, not the $10\times$
ratio, which is set by $P$ and collapses to parity at $P=1$. Two facts
keep the factor in perspective. It multiplies a per-iteration time of
\emph{seconds} of forward-pass inference, not the minutes a backward
pass would take, so the absolute round trip stays inside the modeled
mobile-delay budget \tbd{(the optimized mean per-iteration compute is
${\sim}3.6$\,s, at/under the modeled per-round-trip budget)}. And it
is a per-iteration cost against a per-run benefit: guidance is
designed to reach target accuracy in far fewer iterations, so
run-level compute --- the denominator that sets time-to-accuracy ---
should fall even as the per-iteration cost rises.

% =============================================================================
\section{Ablation}
\label{sec:ablation}

The end-to-end results establish the headline gap and its two causes
at a single operating point. This section justifies the individual
design choices through three angles: \textbf{L1}, enable/disable
studies that show each component actually helps; \textbf{L2},
hyperparameter sensitivity that shows the gain is not brittle; and
\textbf{L3}, design-space alternatives that show the chosen mechanism
beats the options we considered. Qualitative arguments for each choice
live in the design section (\S\ref{sec:design}); here we back them
with data. \tabref{tab:ablation:summary} states which lens applies to
which contribution.

\input{tables/eval/ablation_summary}

We do not ablate async vs.\ sync execution: async is a structural
consequence of dropping the round barrier, so removing it would
re-insert the commitment structure rather than test an independent
mechanism. The value of the iteration-level reframe is instead read
off the round$\to$iteration rung on each selection substrate
independently --- \fwdllm{} $\to$ \fwdllmit{}, \fedbuffp{} $\to$
\fedbuffpit{}, and \felixp{} $\to$ \felixpit{}
(\S\ref{sec:eval:setup}, \S\ref{sec:eval:attribution}) --- each of
which adds per-iteration control while holding that substrate's own
selection policy fixed. This isolates the reframe from both the async
substrate and the C1--C3 contributions, and checking it three times
independently confirms the gain is not an artifact of one particular
selector.

% -----------------------------------------------------------------------
\subsection{L1: does each component help? (enable/disable)}
\label{sec:ablation:ladder}

We build \sys{} up from its guided-JVP-only base, enabling one
mechanism at a time, and measure the effect on peak accuracy and
time-to-accuracy on the same workload ($N{=}100$, $\alpha{=}1$).

\input{figs/code/eval/ablation/abl_optladder}

\textbf{Each enabled mechanism moves peak accuracy up the ladder.}
From the \sys{}-base configuration (guided perturbations only,
$83.0\%$), adding variance-plateau force-commit ($+$C2 stopping) and
then gradient-aware aggregation ($+$C3) raises peak accuracy
monotonically --- $83.0\% \to 82.2\% \to 83.9\% \to 84.1\%$ for the
full stack (\figref{fig:ablation:optladder}). The variance-plateau
step alone slightly trades peak accuracy for earlier commitment; the
gradient-aware aggregation step recovers and exceeds it, and the full
combination is best. This is the enable/disable evidence reviewers ask
for first: each component earns its place.

\begin{plannedexp}{L1 completion: per-contribution on/off}
  \textbf{Goal:} Complete the enable/disable matrix --- guided-off
  (C1), static-$K$/$C$ (C2), naive-weight-agg (C3) --- each toggled
  individually against the full stack, on both availability
  conditions.\\[3pt]
  \textbf{Reuse:} Experiment run-set; no new training beyond the
  toggled configs.\\[3pt]
  \textbf{Expected result:} disabling any single contribution degrades
  time-to-accuracy; C1 dominates the effectiveness lever, C2 the
  utilization lever, C3 the aggregation stability.\\[3pt]
  \textbf{Takeaway:} every contribution is load-bearing; none is
  subsumed by the others.
\end{plannedexp}

% -----------------------------------------------------------------------
\subsection{L3: guided vs.\ cheaper perturbation selection policies}
\label{sec:ablation:c1alt}

\begin{plannedexp}{Guided JVP vs.\ random / cosine / quasi selection}
  \textbf{Goal:} Show guidance is the \emph{only} selection policy
  that moves the needle. The motivation null result
  (\S\ref{sec:motiv:perturbations}) shows all-random, all-cosine, and
  \fwdllm{}'s quasi/alternating policy converge almost identically;
  the ablation re-uses those three cheap policies plus JVP-guided, now
  on accuracy-vs-perturbation-evaluations.\\[3pt]
  \textbf{Baselines:} \sys{} with selection $\in \{$all-random,
  all-cosine, quasi/alternating (\fwdllm{}), JVP-guided$\}$.\\[3pt]
  \textbf{Setup:} 100\% availability (cleanest attribution); fixed
  model, data, $K$/$C$; vary only the selection policy.\\[3pt]
  \textbf{Expected result:} the three cheap policies cluster together;
  JVP-guided reaches target accuracy in the fewest perturbation
  evaluations.\\[3pt]
  \textbf{Takeaway:} cheap alignment scores do not identify directions
  that teach more; JVP guidance is the load-bearing selection signal.
\end{plannedexp}

% -----------------------------------------------------------------------
\subsection{L2: sensitivity to JVP guidance parameters}
\label{sec:ablation:jvp}

\begin{plannedexp}{JVP threshold and refresh frequency sensitivity}
  \textbf{Goal:} Establish that the guided-perturbation gain is not
  brittle to the JVP threshold used to filter directions or to how
  often the server re-runs JVP.\\[3pt]
  \textbf{Baselines:} \sys{} with JVP threshold $\tau \in \{0.1, 0.3,
  0.5\}$ and refresh frequency $\in \{$every iteration, every 10
  iterations, every 50 iterations$\}$.\\[3pt]
  \textbf{Setup:} 100\% availability (cleanest attribution); fixed
  model and task.\\[3pt]
  \textbf{Expected result:} convergence speed is stable across a
  moderate threshold range; very low thresholds approach
  random-sampling performance; very high thresholds reduce direction
  diversity. Refresh every 10 iterations performs comparably to every
  1; every 50 shows small degradation.\\[3pt]
  \textbf{Takeaway:} the JVP hyperparameters can be set without
  per-task tuning.
\end{plannedexp}

\myparagraph{Cost knobs: $P$ and JVP schedule}
The candidate count $P$ and the JVP-computation schedule
(\S\ref{sec:arch:c1}) are the knobs that set the guidance surcharge.
We sweep $P$ (static vs.\ dynamic) to trace the compute/accuracy
frontier, and test the two opposing JVP-schedule hypotheses --- more
JVP early (broad landscape) vs.\ more JVP late (update quality matters
near convergence) --- to fix the default. \tbd{Fill from the $P$-sweep
and JVP-schedule runs.}

% -----------------------------------------------------------------------
\subsection{L1+L2: dynamic vs.\ static $K$/$C$}
\label{sec:ablation:kc}

\begin{plannedexp}{Static vs.\ dynamic $K$/$C$, plus window and factor sensitivity}
  \textbf{Goal:} Show dynamic $K$/$C$ beats the best static setting
  (L1) and is not brittle to the moving-average window $N$ or the
  growth/shrink factors (L2).\\[3pt]
  \textbf{Baselines:} static $K$/$C$ tuned for early phase; static
  tuned for convergence; dynamic $K$/$C$ with $N \in \{5, 10, 20\}$;
  growth/shrink factor $\in \{1.1, 1.25, 1.5\}$; $C_{\max} \in
  \{|$pool $|, 2|$pool$|\}$.\\[3pt]
  \textbf{Setup:} the real-world availability trace; full training run
  (captures both early-phase and convergence behavior).\\[3pt]
  \textbf{Expected result:} dynamic beats both static tunings;
  time-to-accuracy varies $\le$[X]\% across the recommended range;
  aggressive growth/shrink ($1.5\times$) oscillates, conservative
  ($1.1\times$) adapts slowly but still beats static. $N=10$,
  $1.25\times$ recommended default.\\[3pt]
  \textbf{Takeaway:} the $K$/$C$ policy is load-bearing and its
  parameters are not brittle.
\end{plannedexp}

% -----------------------------------------------------------------------
\subsection{L1+L3: gradient-aware vs.\ naive aggregation}
\label{sec:ablation:c3}

\begin{plannedexp}{Direction-aware aggregation vs.\ naive weight averaging}
  \textbf{Goal:} Show \sys{}'s direction-aware, scalar-decay
  aggregation (C3, \S\ref{sec:arch:c6}) beats naive averaging and a
  single-scalar staleness discount, and test its knobs.\\[3pt]
  \textbf{Baselines:} naive average; single FeLiX-style (TierFuse)
  staleness$\times$utility discount~\cite{garg2026felix}; \sys{} C3
  with $\gamma$ sweep, greedy vs.\ cluster novelty, hard vs.\ soft
  sign gate.\\[3pt]
  \textbf{Setup:} the real-world availability trace; fixed model and
  task.\\[3pt]
  \textbf{Expected result:} the single-scalar discount wrongly
  penalizes the fixed direction $v$ and underperforms; C3's
  direction/scalar split is more stable under staleness.\\[3pt]
  \textbf{Takeaway:} forward-gradient updates need direction-aware
  aggregation; a borrowed weight-merge discount is wrong in kind.
  \tbd{Gated on C3 validation and formula sync with the latest code
  (\S\ref{sec:arch:c6}).}
\end{plannedexp}

% -----------------------------------------------------------------------
\subsection{L2+L3: databin size}
\label{sec:ablation:databin}

\begin{plannedexp}{Databin-size sweep}
  \textbf{Goal:} Locate the databin size that balances the two failure
  modes named in \S\ref{sec:arch:global} --- small bins
  over-communicate, large bins wash out per-client signal.\\[3pt]
  \textbf{Baselines:} \sys{} across a databin-size sweep (small $\to$
  large), fixed model, data, and $K$/$C$ policy.\\[3pt]
  \textbf{Setup:} 100\% availability for a clean read on the
  communication/learning trade, then confirm on the real-world
  trace.\\[3pt]
  \textbf{Expected result:} communication falls monotonically with
  databin size while learning-per-databin peaks at an intermediate
  size; time-to-accuracy is U-shaped, with the minimum away from both
  extremes.\\[3pt]
  \textbf{Takeaway:} databin size is a genuine trade, and the default
  sits between the over-communication and signal-washout regimes.
\end{plannedexp}

% -----------------------------------------------------------------------
\subsection{Robustness under non-IID data distribution}
\label{sec:ablation:noniid}

\begin{plannedexp}{Non-IID severity: Dirichlet $\alpha$ sweep}
  \textbf{Goal:} Confirm that \sys{}'s gain over \fwdllm{} holds
  across non-IID severity levels.\\[3pt]
  \textbf{Baselines:} \fwdllm{}, \fwdllmito{}, and \sys{} at $\alpha
  \in \{0.5$ (moderate non-IID$), 1.0$ (mild, the primary
  condition$)\}$; \tbd{extend to \fedbuffp{}/\felixp{} once their runs
  land}. The more severe $\alpha{=}0.1$ is excluded (learning too slow
  across all baselines to complete convergence runs).\\[3pt]
  \textbf{Setup:} the real-world availability trace. Fixed model and
  task.\\[3pt]
  \textbf{Expected result:} \sys{} maintains its speedup at both
  $\alpha$ values; the \emph{relative} gain is stable or improves,
  since guided perturbations select locally informative directions
  regardless of data distribution.\\[3pt]
  \textbf{Takeaway:} \sys{}'s gain is robust across the tested range
  of non-IID severity; the residual non-IID bias noted in
  \S\ref{sec:background} does not eliminate it.
\end{plannedexp}

% =============================================================================
% Preempts not addressed inline
% =============================================================================

\preempt
  {Results on DistilBERT do not generalize across model families.}
  {here} {Evaluate on $\ge$2 further families (LLaMA2-7B primary;
  Mistral-7B secondary). If only one is available by the deadline,
  scope the generalizability claim explicitly and commit to the second
  family in the submission.}

\preempt
  {\sys{} underperforms \fwdllm{} in some easy condition (potential
  regression).} {here} {If dynamic $K$/$C$ or JVP overhead adds
  latency that slows the very first training phase before guidance
  kicks in: report the regression, show it is transient (bounded to
  $\le$[X] iterations), and argue acceptability in light of the
  full-run and low-availability gains. Do not suppress.}

\preempt
  {SPRY is more accurate AND more robust to availability --- \sys{}'s
  value is unclear.} {here} {If SPRY without availability-aware
  control degrades similarly to \fwdllm{} at low availability
  (expected), then \sys{}'s deployment-model advantage (no weight
  splitting) is the primary differentiator, not gradient quality.
  Frame accordingly.}
